\scanpage{123}\mspage{62}

de notre relation d'équivalence. $[\![$ C'est le moment peut-être, d'ailleurs, de noter la relation plausible

\[
Z'\text{ est un groupe \uline{commutatif}}
\tag*{$(48)$}
\]

% slanted
\margin{Je ne vais pas la poser en axiome pour \uline{l'instant} cf la démonstration bien \unsure{foireuse|foirée} suivante}

Car si $\lambda,\mu\in Z'$, et $(s,t)\in S^{\unsure{2}}$, on a

\diagram{\small\begin{tikzcd}[ampersand replacement=\&,column sep=small,row sep=large]
(s,t)\in A_0 \arrow[r,Leftrightarrow,"\lambda\in Z'"] \arrow[d,Leftrightarrow,"\lambda\mu\in Z'"'] \&
 \{\lambda s,\lambda^{-1}t\}\in A_0 \arrow[r,Leftrightarrow,"\mu\in Z'"] \&
 \{\mu\lambda s,\mu^{-1}\lambda^{-1}t\}\in A_0\\
 \{(\lambda\mu)s,(\lambda\mu)^{-1}t\}\in A_0\mathpunct{,} \& \& 
\end{tikzcd}}

donc en comparant les extrêmes, et posant $\mu^{-1}\lambda^{-1}t=(\lambda\mu)^{-1}t=u$

\[
\{u,(\lambda\mu)s\}\in A_0\Longleftrightarrow\{u,(\mu\lambda)s\}\in A_0
\]

$\forall s,u\in S$, ce qui s'écrit \deleted{\unsure{ce}}

\[
\deleted{\unsure{(49)}}\quad A_0((\lambda\mu)s)=A_0((\mu\lambda)(s))\ ,
\]

égalité des \uline{étoiles} de $(\lambda\mu)(s)$ et $(\mu\lambda)(s)$.

\margin{\deleted{\unsure{g}}}

Si donc on suppose que l'un sommet est déterminé par son étoile, on trouve bien (48) -- mais je ne vois pas \deleted{faire} reposer l'axiome explicite sur cette détermination un sommet par son étoile -- (qui est d'ailleurs \deleted{faux} pour l'octaèdre\ldots). $]\!]$

% slanted
\margin{c'est vrai du fait que $P_0$ est simpl. triangulé \unsure{connexe|convexe} dont \uline{tt} les sommets sont de \unsure{égal|même} ordre, à l'exception du seul octaèdre}

% slanted
\margin{cf bas de la page 22 pour un argument bizarre \unsure{en ce sens|réussi|aussi}\ldots}
% archive sheet number: 122
